\documentclass[10pt,a4paper]{amsart}
\usepackage[margin=19mm]{geometry}
\usepackage{amsmath,amssymb,mathtools}
\usepackage[scr=rsfs,cal=euler]{mathalfa}
\usepackage{xfrac}
\usepackage[unicode,hidelinks,bookmarks=false]{hyperref}
\newcommand{\ie}{i.e.\ }
\newcommand{\cf}{cf.\ }
\newcommand{\resp}{respectively\ }
\newcommand{\Subsection}{\subsection}
\providecommand{\nobreakdash}{\nobreak\hbox{-}}
\setlength{\emergencystretch}{2em}
\multlinegap0pt
\usepackage{float}
\usepackage[shortlabels]{enumitem}
\usepackage{amssymb}
\usepackage{mathtools}
\usepackage{bm}
\newtheorem{theorem}{Theorem}
\newcommand{\Z}{\mathbb{Z}}
\newcommand{\lm}{\textup{\texttt{-}}}
\newcommand{\lp}{\textup{\texttt{+}}}
\title{Atoms meet symbols}
\author{Leonardo F. Cavenaghi, Ludmil Katzarkov,, Maxim Kontsevich}
\date{Source: arXiv:2509.15831v4 (CC BY 4.0)}
\begin{document}
\maketitle

\noindent\emph{Source and licence.} Atoms meet symbols. Version arXiv:2509.15831v4 (CC BY 4.0), distributed by arXiv under the Creative Commons Attribution 4.0 International licence, http://creativecommons.org/licenses/by/4.0/, as recorded in the arXiv licence metadata for this exact version and shown on the abstract page https://arxiv.org/abs/2509.15831v4. Full licence notice: \texttt{licenses/arxiv-2509.15831-CC-BY-4.0.txt}.\par\smallskip

\noindent\textit{Editorial scope.} Counted unit: the article's theorem thm:pure-geometry (source0.tex lines 2345--2379). The article's complete original proof of it is delivered with it (source0.tex lines 2382--2510, one \texttt{proof} environment of 129 lines). No mathematics is written, repaired or re-derived: the statement and the proof are the article's own lines, in the article's own notation, and no retained line is substituted or re-flowed.  No further source-local context is needed: every cross-reference of every retained line resolves inside the two delivered spans.  Nothing was omitted from any retained span, so the package carries no omission record.  No retained line contains a citation, so the wrapper carries no bibliography and no bibliography entry.  No retained line contains a figure environment, an image-inclusion command or drawing code, so no image is delivered and the package declares no asset dependency.  Editorial formatting only, in addition: an AMS \texttt{amsart} preamble that restates the article's own declarations and macros used by the retained lines, namely the environments \texttt{theorem} on separate counters, together with the standard packages those lines need.  The retained environments print in this document as Theorem 1 in order of appearance, whereas the article prints the same items with the numbering of its own full document.  The complete native source of the article as distributed in the arXiv e-print for this exact version is bundled unchanged as \texttt{sources/185327/source0.tex}.
\begin{theorem}\label{thm:pure-geometry}
Let $X$ be a smooth projective threefold with a generically free $\mathbb Z/3$‐action.  Then the following quantity is invariant under \(\mathbb{Z}/3\)-equivariant blowups:
\begin{equation}
  J = 
    \sum_{[\lp\lp\lm],\; [\lp\lm\lm]} 1
    + 
    \sum_{(C,[\lp\lp]),\; (C,[\lm\lm])} \left(1 - \mathsf{g} + d\right)
    +
    \sum_{(C,[\lp\lm])} \left(2 - 2\mathsf{g} + d\right)
    +
    \sum_{S} \left(3 - 3\mathsf{g} - K_S\cdot \mathcal N\right)
\end{equation}
where the summation runs over fixed point set components $X^{\mathbb Z/3}$ where:
\begin{itemize}
  \item 
   $C$ is a curve with genus \(\mathsf{g}=\mathsf{g}(C) \geq 0\).
  \item 
    \(d \in \mathbb{Z}\) is the degree of \(\wedge^{2}(\mathcal{N}_XC)\), where \(\mathcal{N}_XC\) is the normal bundle of $C$ in $X$.
  \item $S$ is birational to a ruled surface whose base of the ruling is a curve with genus 
    \(\mathsf{g} \geq 0\) (slightly abusing the notation, we write $\mathsf g=\mathsf g(S)$).
  \item 
    \(K_{S}\) is the canonical bundle of the surface \(S\).
  \item 
    \(\mathcal{N}=\mathcal N_X S\) is the normal bundle of \(S\) in $X$.
\end{itemize}
The signs decorations are the non-zero characters for \(\mathbb Z/3 = (\mathbb Z/3)^{\vee}\), i.e.,
\[
  \lp,\,\lm \;=\; 1,\,2 \pmod{3}.
\]
so that
\begin{itemize}
    \item $[\lp\lp\lm], [\lp\lm\lm]$ are the normal weights for the $\mathbb Z/3$-normal bundle representation at fixed points
    \item for symbols like $(C,[\lp\lp]), (C,[\lp\lm]), (C,[\lm\lm])$ the terms $[\lp\lp], [\lm\lm], [\lp\lm]$ stand to the normal weights for the $\mathbb Z/3$-normal bundle representation at the fixed curve $C$.
\end{itemize}
 \end{theorem}

\begin{proof}[Proof of Theorem \ref{thm:pure-geometry}]
To establish the invariance of the quantity $J$, we must show that for any smooth, invariant, closed center of blowup $Z \subset X$ (of codimension $r \geq 2$), the change $\Delta J:= J(\mathrm{Bl}_Z X) - J(X)$ is zero. We proceed by a case-by-case analysis of all possible centers $Z$. The locus of fixed points on the blowup maps to the locus of fixed points on the original variety $X$. We separate contributions by connected components of $X^{\Z/3}$ of dimensions $0,1,2$, and the local behavior of the center $Z$ of the blowup near these components. 

\noindent \textbf{Case 1: Blowup in an isolated fixed point.}

\begin{itemize}
    \item \textbf{Type $[\lp\lp\lp]$} (\textit{and the opposite case  $[\lm\lm\lm]$}): A fixed point of this type contributes $0$ to $J$. The blowup replaces the point with an exceptional divisor $E \cong \mathbb{P}^2$, which is a new surface component in the $\mathbb Z/3$-fixed locus on the blowup. For this new surface $S=E$, we have genus $\mathsf{g}(S)=0$, canonical bundle $K_S = \mathcal{O}_{\mathbb{P}^2}(-3)$, and normal bundle $\mathcal{N}=\mathcal{O}_{\mathbb{P}^2}(-1)$. The contribution of this new surface to $J$ is:
    \begin{align}
        3 - 3\mathsf{g}(S) - K_S \cdot \mathcal{N} = 3 - 3(0) - (-3)(-1) = 3 - 3 = 0.
    \end{align}
    Thus, $\Delta J = \underbrace{0}_{\text{new surface}} - \underbrace{0}_{\text{old point}} = 0$.

    \item \textbf{Type $[\lp\lp\lm]$} (\textit{and the opposite case}): A fixed point of type $[\lp\lp\lm]$ contributes $1$ to $J$. The blowup replaces the point with an exceptional divisor $E$. The $\mathbb Z/3$-fixed locus on $E$ is a curve $C \cong \mathbb{P}^1$ with normal weights $[\lp\lp]$. For this curve, $\mathsf{g}(C) = 0$ and the degree of the second exterior power of its normal bundle is $d = \deg(\wedge^2 \mathcal{N}_{X'}C) = 0$. Its contribution to $J$ is:
    \begin{align}
        1 - \mathsf{g}(C) + d = 1 - 0 + 0 = 1.
    \end{align}
    Thus, $\Delta J = \underbrace{1}_{\text{new curve}} - \underbrace{1}_{\text{old point}} = 0$.
\end{itemize}

\noindent \textbf{Case 2: Blowup in a curve passing through an isolated fixed point.}

\begin{itemize}
    \item \textbf{Types $[\lp\lp\lp]$} (\textit{and the opposite case}):  The initial fixed point contributes $0$. Blowing up in an invariant curve passing through it replaces the point with a fixed curve $C \cong \mathbb{P}^1$ with normal weights $[\lp\lp]$. The degree of the second exterior power of its normal bundle is $d=-1$. The new contribution is:
    \begin{align}
        1 - \mathsf{g}(C) + d = 1 - 0 + (-1) = 0.
    \end{align}
    Hence, $\Delta J = \underbrace{0}_{\text{new curve}} - \underbrace{0}_{\text{old point}} = 0$.

    \item \textbf{Type $[\lp\lp\lm]$ and blowup in the $\lm$ direction}(\textit{and the opposite case}): The initial point contributes $1$. The blowup replaces it with a new fixed curve $C \cong \mathbb{P}^1$ with normal weights $[\lp\lm]$ and total degree $d=d_\lp+d_\lm = 0+(-1)=-1$. The new contribution is:
    \begin{align}
        2 - 2\mathsf{g}(C) + d = 2 - 2(0) + (-1) = 1.
    \end{align}
    Hence, $\Delta J =\underbrace{1} _{\text{new curve}} - \underbrace{1}_{\text{old point}} = 0$.

    \item \textbf{Type $[\lp\lp\lm]$ and blow up in the $\lp$ direction}(\textit{and the opposite case}):
     The initial point contributes $1$. The blowup replaces it by two isolated fixed points of types $[\lp\lp\lp]$ and $[\lp\lm\lm]$. The new contribution is $1$, hence $\Delta J=\underbrace{1}_{\text{new point of type }[\lp\lm\lm]}-\underbrace{1}_{\text{old point}}=0$.
\end{itemize}

\noindent \textbf{Case 3: Blowup in a point on a fixed curve.}

\begin{itemize}
    \item \textbf{The curve of type $[\lp\lp]$}(\textit{and the opposite case}):  Let the original curve be $C$ with contribution $2 - 2\mathsf{g} + d$. Blowing up at a point on $C$ produces in the fixed locus a new component $\mathbb P^1$ with normal weights [\lp\lm] and the degree of the normal bundle equal to $-1+1=0$. The curve $C$ is replaced by its proper transform $C'$, for which the degree of the second exterior power of the normal bundle becomes $d' = d-2$. The total change is:
    \begin{align}
        \Delta J &= \underbrace{2-2\cdot 0+0}_{\text{new curve }\mathbb P^1} + \underbrace{(2 - 2\mathsf{g} + (d-2))}_{\text{new curve}} - \underbrace{(2 - 2\mathsf{g} + d)}_{\text{old curve}} \\
        &= 2 + (2 - 2\mathsf{g} + (d - 2)) - (2 - 2\mathsf{g} + d) = 0.
    \end{align}
    \item \textbf{The curve of type $[\lp\lm]$:} Let the original curve be $C$ with contribution $2 - 2\mathsf{g} + d$. Blowing up a fixed point on $C$ replaces that point with two new isolated fixed points of types $[\lp\lp\lm]$ and $[\lp\lm\lm]$, each contributing $1$. The curve $C$ is replaced by its proper transform $C'$, for which the degree of the second exterior power of the normal bundle becomes $d' = d-2$. The total change is:
    \begin{align}
        \Delta J &= \underbrace{(1 + 1)}_{\text{new points}} + \underbrace{(2 - 2\mathsf{g} + (d-2))}_{\text{new curve}} - \underbrace{(2 - 2\mathsf{g} + d)}_{\text{old curve}} \\
        &= 2 + (2 - 2\mathsf{g} + (d - 2))
        - (2 - 2\mathsf{g} + d) = 0.
    \end{align}
\end{itemize}

\noindent \textbf{Case 4: Blowup in a fixed curve.}
\begin{itemize}
    \item \textbf{The curve of type $[\lp\lp]$} (\textit{and the opposite case}): The curve $C$ contributes $1 - \mathsf{g} + d$. It is replaced by a ruled surface $S = \mathbb{P}(\mathcal{N}_XC)$ fibered over $C$, with $\mathsf{g}(S) = \mathsf{g}(C) = \mathsf{g}$. The new contribution is $3 - 3\mathsf{g} - K_S \cdot \mathcal{N}$ where $\mathcal N=\mathcal N_S \widetilde X$ is the normal line bundle to $S$ in the blowup $\widetilde X=\mathrm{Bl}_C X$. In order to simplify the calculation, we assume that the rank two bundle $\mathcal N_X C$ splits into the sum $\mathcal L_1\oplus \mathcal L_2$ of line bundles of certain degrees $d_1,d_2$ such that $d=d_1+d_2$. The general case of a non-split bundle follows by continuity. To calculate $K_S \cdot \mathcal{N}$, we use the basis of $\mathsf{NS}(S)$ given by the curve $C_1 = \mathbb{P}(\mathcal{L}_1)\subset S$ and the fiber $C_{\text{vert}} \cong \mathbb{P}^1$. The intersection numbers with $K_S$ and $\mathcal{N}$ are known:
    \begin{align}
        K_S \cdot C_1 &= 2\mathsf{g}-2 + d_1-d_2,  & K_S \cdot C_{\text{vert}} &= -2 \\
        \mathcal{N} \cdot C_1 &= d_1, & \mathcal{N} \cdot C_{\text{vert}} &= -1
    \end{align}
    The intersection form on $\mathsf{NS}(S)$ in the basis $(C_1,C_{\mathrm{vert}})$ is given by
    $$\begin{pmatrix}d_2-d_1 & 1\\1& 0\end{pmatrix}$$
    Therefore, the intersection $K_S\cdot \mathcal N$ is given by
    $$ \begin{pmatrix}d_1  &-1\end{pmatrix} \cdot \begin{pmatrix}d_2-d_1 & 1\\1& 0\end{pmatrix}^{-1}\cdot  \begin{pmatrix}2\mathsf{g}-2+d_1-d_2\\-2\end{pmatrix} =2-2\mathsf{g}-d_1-d_2.  $$
    The change in $J$ is therefore:
    \begin{align}
       \Delta J &=\underbrace{ (3 - 3\mathsf{g} - (2 - 2\mathsf{g} - d))}_{\text{new surface}} - \underbrace{(1 - \mathsf{g} + d)}_{\text{old curve}} \\
        &= (1 - \mathsf{g} + d) - (1 - \mathsf{g} + d) = 0.
    \end{align}

    \item \textbf{The curve of type $[\lp\lm]$:} The curve $C$ contributes $2 - 2\mathsf{g} + d$. It is replaced by two curves, $C_1$ (type $[\lp\lp]$) and $C_2$ (type $[\lm\lm]$). Let denote the degrees of $\lp$ and $\lm$ components of $\mathcal N_XC$ by $d_\lp$ and $d_\lm$ respectively, so $d = d_\lp + d_\lm$. The total change is:
    \begin{align}
        \Delta J &= \underbrace{(1 - \mathsf{g} + d_\lp)}_{\text{new curve } C_1} + \underbrace{(1 - \mathsf{g} + d_\lm)}_{\text{new curve } C_2} - \underbrace{(2 - 2\mathsf{g} + d)}_{\text{old curve } C} \\
        &= (2 - 2\mathsf{g} + d_\lp + d_\lm) - (2 - 2\mathsf{g} + d_\lp + d_\lm) = 0.
    \end{align}
\end{itemize}

\noindent \textbf{Case 5: Blowup in a curve transversal to a fixed curve.}
\begin{itemize}
    \item \textbf{The fixed curve of type $[\lp\lp]$} (\textit{and the opposite case}):  The original curve $C$ contributes $1-\mathsf{g}+d$. For each intersection point the blowup adds a new isolated fixed point of type $[\lp\lp\lm]$. The proper transform of $C$, denoted $C'$, is still a component of the fixed locus. The degree $d'$ of the normal bundle of $C'$ is equal to $d'=d-k$, where $k$ is the number of intersection points of $C$ and the center of the blowup. The total change is:
 \begin{align}
    \Delta J &= k (\text{contribution of new points}) \\
             &\quad + (1 - \mathsf g + d - k) (\text{contribution of } C') \\
             &\quad - (1 - \mathsf g + d) (\text{contribution of } C) = 0
\end{align}
\end{itemize}

\noindent \textbf{Case 6: Blowup in a point on a fixed surface.}

\begin{itemize}
    \item \textbf{The fixed surface of type $\lp$} (\textit{and the opposite case}):  The surface $S$ contributes $3 - 3\mathsf{g} - K_S \cdot \mathcal{N}$. Blowing up a point on $S$ creates a new isolated fixed point of type $[\lp\lm\lm]$ and replaces $S$ with its blowup $\widetilde{S}$. For $\widetilde{S}$, we have $K_{\widetilde{S}} = \pi^{\ast}K_S + E$ and $\widetilde{\mathcal{N}} = \pi^{\ast}\mathcal{N} - E$, where $E$ is the exceptional divisor. The new intersection product is $K_{\widetilde{S}} \cdot \widetilde{\mathcal{N}} = (\pi^{\ast}K_S + E) \cdot (\pi^{\ast}\mathcal{N} - E) = K_S \cdot \mathcal{N} - E^2 = K_S \cdot \mathcal{N} + 1$. The change in $J$ is:
    \begin{align}
        \Delta J &= \underbrace{1}_{\text{new point}} + \underbrace{(3 - 3\mathsf{g} - (K_S\cdot\mathcal{N} + 1))}_{\text{new surface}} - \underbrace{(3 - 3\mathsf{g} - K_S\cdot\mathcal{N})}_{\text{old surface}} \\
        &= 1 - 1 = 0.
    \end{align}
\end{itemize}

\noindent \textbf{Case 7: Blowup in a transverse curve to a fixed surface.}
\begin{itemize}
    \item  \textbf{The fixed surface of type $\lp$} (\textit{and the opposite case}): Let $S$ be a fixed surface and $C$ be a $\Z/3$-invariant curve that intersects $S$ transversely at a finite set of points.  The blowup of $X$ in $C$ replaces component $S$ of the fixed locus with its proper transform $\widetilde{S}$, which is isomorphic to the blowup of $S$ at the points $C \cap S$. The canonical and normal bundles transform as $K_{\widetilde{S}} = \pi^{\ast}K_S + \sum E_i$ and $\widetilde{\mathcal{N}} = \pi^{\ast}\mathcal{N}$, where $E_i$ are the exceptional divisors on $\widetilde{S}$. Since $\pi^{\ast}\mathcal{N} \cdot E_i = 0$, the intersection product $K_{\widetilde{S}}\cdot\widetilde{\mathcal{N}} = K_S\cdot\mathcal{N}$ is unchanged. 
    
    The change in $J$ is therefore $$\Delta J = \underbrace{(3-3\mathsf g +K_{\widetilde S}\cdot \widetilde{\mathcal N} )}_{\text{new surface}} -   (\underbrace{3-3\mathsf g +K_{S}\cdot \mathcal N)}_{\text{old surface}}= 0.$$
\end{itemize}

\noindent \textbf{Case 8: Blowup in a curve on a fixed surface.}

\begin{itemize}
    \item  \textbf{The fixed surface of type $\lp$} (\textit{and the opposite case}): Let $C \subset S$ be a closed curve. In the blowup $\mathrm{Bl}_CX$, the component $S$ of the fixed locus is replaced by its proper transform $\widetilde{S}\cong S$, and a new curve $\widetilde C\cong C$, with normal weights $[\lp\lm]$. The change in the surface term's contribution is $K_S \cdot C$.
    \begin{align}
        \Delta(3-3\mathsf{g}-K_S\cdot\mathcal{N}) &= -(K_{\widetilde{S}}\cdot\widetilde{\mathcal{N}}) - (-K_S\cdot\mathcal{N})\\
        &= -K_S\cdot(\mathcal{N}-C) + K_S\cdot\mathcal{N}\\
        &= K_S \cdot C= 2\mathsf{g}(C) - 2 - \mathcal N_S C\cdot C.
    \end{align}
    The last equality is the adjunction formula. Now turn to the new curve component  $\widetilde C$ of the fixed locus. Its normal bundle canonically splits into $\lp$ and $\lm$ parts. The degrees these line bundles are $d_\lp = \mathcal{N}_XS \cdot C$ and $d_\lm = \mathcal N_S C\cdot C - \mathcal{N}_XS \cdot C$ respectively. The contribution of this new curve to $J$ is therefore:
    \begin{align}
        J_{\text{new curve}} &= 2 - 2\mathsf{g}(C') + d_\lp + d_\lm \\
        &= 2 - 2\mathsf{g}(C) + (\mathcal{N}_XS \cdot C) + (\mathcal N_S C\cdot C - \mathcal{N}_XS \cdot C) \\
        &= 2 - 2\mathsf{g}(C) + \mathcal N_S C\cdot C.
    \end{align}
    The total change in $J$ is:
    \begin{align}
        \Delta J &= J_{\widetilde S}-J_S+J_{\text{new curve}}\\
        &= (2\mathsf{g}(C) - 2 -  \mathcal{N}_XS \cdot C) + (2 - 2\mathsf{g}(C) +  \mathcal{N}_XS \cdot C) =0.
    \end{align}
\end{itemize}

Since $\Delta J = 0$ in all cases, the quantity $J$ is invariant under $\mathbb{Z}/3$-equivariant blowups.
\end{proof}

\end{document}
